\documentclass[a4paper,12pt]{article}
\usepackage[utf8]{inputenc}
\usepackage[T1]{fontenc}
\usepackage[french]{babel}
\usepackage[right=2.5cm, left=2.5cm]{geometry}
\usepackage[ddmmyyyy]{datetime}
\usepackage[table]{xcolor}
\usepackage{lmodern,mathptmx,changepage,titlesec,hyperref,listings,lstautogobble,graphicx,array,longtable,multirow,lipsum,tikz,shorttoc,enumitem}
\usetikzlibrary{arrows,automata}
\usetikzlibrary{positioning}

\renewcommand{\rmdefault}{\sfdefault} %Utilisation de la police sans-serif ("Computer Modern Sans") pour la police roman
\renewcommand{\ttdefault}{pcr} 	%Utilisation d'une police "CourrierNew" pour la police monospaced (pour faire un listing manuel)
\linespread{1.15}				%Interligne

%Utilisation de liens colorés en bleu et soulignés
\hypersetup{colorlinks=true, urlcolor=blue, urlbordercolor=blue, linkcolor=black, linkbordercolor=white}
\makeatletter \Hy@AtBeginDocument{\def\@pdfborder{0 0 1} \def\@pdfborderstyle{/S/U/W 1}}\makeatother

\titlespacing*{\section} {0cm}{7ex plus 1ex minus .2ex}{1.5ex plus .2ex}
\titlespacing*{\subsection} {0cm}{4.5ex plus 1ex minus .2ex}{1.5ex plus .2ex}
\titleformat*{\section}{\huge\bfseries}
\titleformat*{\subsection}{\Large\bfseries}
\titleformat*{\subsubsection}{\normalsize\bfseries}

\definecolor{darkgreen}{rgb}{0,0.8,0}
\definecolor{mygray}{rgb}{0.93,0.93,0.93}
\definecolor{mymauve}{rgb}{0.58,0,0.82}
\lstset{	
	basicstyle=\small\ttfamily,
	backgroundcolor=\color{mygray},
	breaklines=true,
	breakatwhitespace=true,
	postbreak=\raisebox{0ex}[0ex][0ex]{\ensuremath{\color{red}\hookrightarrow\space}},
	tabsize=3,
	frame=none,
	rulecolor=\color{black},
	keywordstyle=\color{blue}\bfseries,
	stringstyle=\color{orange},
	showstringspaces=false,
	commentstyle=\footnotesize\color{darkgreen},
	keepspaces=true,
	extendedchars=true,
	numbers=left,
	numberstyle=\tiny\color{lightgray},
	stepnumber=1,
	escapeinside={(@}{@)},
	autogobble=true,
	literate=
		{á}{{\'a}}1 {é}{{\'e}}1 {í}{{}}1 {ó}{{\'o}}1 {ú}{{\'u}}1
		{Á}{{\'A}}1 {É}{{\'E}}1 {Í}{{\'I}}1 {Ó}{{\'O}}1 {Ú}{{\'U}}1
		{à}{{\`a}}1 {è}{{\`e}}1 {ì}{{\`i}}1 {ò}{{\`o}}1 {ù}{{\`u}}1
		{À}{{\`A}}1 {È}{{\'E}}1 {Ì}{{\`I}}1 {Ò}{{\`O}}1 {Ù}{{\`U}}1
		{ä}{{\"a}}1 {ë}{{\"e}}1 {ï}{{\"i}}1 {ö}{{\"o}}1 {ü}{{\"u}}1
		{Ä}{{\"A}}1 {Ë}{{\"E}}1 {Ï}{{\"I}}1 {Ö}{{\"O}}1 {Ü}{{\"U}}1
		{â}{{\^a}}1 {ê}{{\^e}}1 {î}{{\^i}}1 {ô}{{\^o}}1 {û}{{\^u}}1
		{Â}{{\^A}}1 {Ê}{{\^E}}1 {Î}{{\^I}}1 {Ô}{{\^O}}1 {Û}{{\^U}}1
		{œ}{{\oe}}1 {Œ}{{\OE}}1 {æ}{{\ae}}1 {Æ}{{\AE}}1 {ß}{{\ss}}1
		{ç}{{\c c}}1 {Ç}{{\c C}}1 {ø}{{\o}}1 {å}{{\r a}}1 {Å}{{\r A}}1
		{€}{{e}}1 {£}{{\pounds}}1 {«}{{\guillemotleft}}1
		{»}{{\guillemotright}}1 {ñ}{{\~n}}1 {Ñ}{{\~N}}1 {¿}{{?`}}1
}

%Redéfinition de la taille de \Huge pour le titre du document
\makeatletter\renewcommand\Huge{\@setfontsize\Huge{37pt}{40}}\makeatother
\date{\today}

\title{\vspace{\fill}\textbf{\Huge Compte Rendu}}
\author{
	Sonny Klotz - Jean-Didier Pailleux - Malek Zemni
	\vspace{2em}\\
	\textit{Interface de chargement, de contrôle}\\\textit{et d’analyse statistique des données}\\\textit{pour la constitution d’un graphe de flux}
	\vspace{2em}
}

\begin{document}
\pagenumbering{gobble}\clearpage
\maketitle\vspace{9em}
\begin{center}\includegraphics[scale=0.7]{../Cahier/logo.png}\end{center}
\begin{flushright}Module \textit{Projet}\end{flushright}
\newpage
\tableofcontents
\newpage\clearpage\pagenumbering{arabic}

	\section*{Introduction}
	Ce document est le compte-rendu final de notre travail qui s'inscrit dans le cadre du module \textit{Projet} de la licence informatique de l' \textit{\textbf{UVSQ}}. Le sujet de ce projet a été proposé par l'entreprise \textit{\textbf{DCbrain}}.\\
	
	Notre client, \textit{\textbf{DCbrain}}, développe des outils qui permettent de visualiser le comportement des \textbf{réseaux physiques} afin de prédire et trouver les problèmes de ces réseaux dans le but de les optimiser. Ces réseaux physiques sont principalement les réseaux industriels, de fluide ou de distribution, tels les réseaux électriques. Des données sont collectées à partir de ces réseaux puis analysées grâce aux technologies du \textbf{Big Data}. Le but est d'en extraire une valeur ajoutée.\\
	
	Notre tâche correspond au travail indispensable d’analyse préliminaire de ces données, afin de faciliter leur exploitation. Ce projet a donc pour objectif de fournir aux utilisateurs une application web qui permettra de charger les données collectées, les visualiser et enfin les analyser.\\
	
	Dans une première partie de ce document, on présentera l'architecture de notre application, illustrée par un organigramme qui a été préalablement établi. Cet organigramme représente un découpage de l'application en modules avec les différentes informations qui circulent entre ces modules.\\
	Dans une deuxième partie, on parlera des différents langages de programmation choisis pour réaliser l'application et des contraintes qui ont justifié ces choix.\\
	Ensuite, dans une troisième partie, on traitera la partie technique de notre projet, c'est à dire, l'application, son fonctionnement et ses problèmes.\\
	Finalement, dans les deux dernières parties, on établira un bilan quant à l'organisation interne au sein du groupe et un bilan comparatif des coûts présumés et des coûts finaux de notre produit.\\
	
	%PAS ICI 
	%Deux documents supplémentaires sont livrés avec l'application, un manuel utilisateur décrivant comment faire fonctionner notre application, et une documentation technique des modules de notre produit.
			
	%Peut-être parler un peu plus du big data, aussi des techniques de graphes mining dans le travail de DC brain, pour présenter le fait que le réseau est assimilé à un graphe =============== oh no c bcp
	
	\section{Architecture de l'application}
		\subsection{Organigramme et données échangées}
		Cet organigramme a été préalablement établi dans le cahier des charges du produit. Il représente la décomposition en modules de l'application, ainsi que les informations qui circulent entre ces modules :
		\begin{figure}[H]
			\begin{tikzpicture}
			\begin{scope}[xscale=2,yscale=1.5]	
				\node (IW) at (0,5) [rectangle,draw,text depth=3cm,minimum width=16cm,minimum height=4cm,font=\textbf\Large] {\begin{tabular}{c}Interface web\end{tabular}};
				\node (F2) [rectangle,draw,dashed] at ([yshift=0.4cm]IW.center) {\begin{tabular}{c}Fenêtre rôle et choix colonne \end{tabular}};
				\node (F1) [rectangle,draw,dashed] at ([xshift=-4cm]F2.east) {\begin{tabular}{c}Fenêtre choix fichier\end{tabular}};
				\node (F3) [rectangle,draw,dashed] at ([xshift=4.1cm]F2.west) {\begin{tabular}{c}Fenêtre résultats ADD\end{tabular}};
				\node (MAIN) [rectangle,draw,dashed,below=of F2.south,yshift=0cm] {\begin{tabular}{c}Gestion des flux\end{tabular}};

			
				\node (API1) at (-2.5,0) [rectangle,draw,fill=blue!25,text depth=-3.5cm,minimum width=5cm,minimum height=5cm,font=\textbf\Large] {\begin{tabular}{c}API 1\\Chargement des données\end{tabular}};
				\node (VERIF) [rectangle,draw,dashed,fill=white] at ([yshift=1cm]API1.center){\begin{tabular}{c}Vérification format\\fichier\end{tabular}};
				\node (ANALYS) [rectangle,draw,dashed,fill=white,below=of VERIF.south,yshift=0.5cm] {\begin{tabular}{c}Analyse contenu\\fichier\end{tabular}};
			
				\node (API2) at (2.5,0) [rectangle,draw,fill=blue!25,text depth=-5cm,minimum width=5cm,minimum height=6.5cm,font=\textbf\Large] {\begin{tabular}{c}API 2\\Analyse descriptive des données\end{tabular}};
				\node (ADD1) [rectangle,draw,dashed,fill=white] at ([yshift=1.7cm]API2.center){\begin{tabular}{c}ADD qualitatives\end{tabular}};
				\node (ADD2) [rectangle,draw,dashed,fill=white,below=of ADD1.south,yshift=0.5cm] {\begin{tabular}{c}ADD quantitatives\\discrètes\end{tabular}};
				\node (ADD3) [rectangle,draw,dashed,fill=white,below=of ADD2.south,yshift=0.5cm] {\begin{tabular}{c}ADD quantitatives\\continues\end{tabular}};
			
				\draw[-triangle 45] (F1.south east) -- node[anchor=south,left]{11} (MAIN.north west);
				\path[->,>=stealth'] (MAIN) edge[bend left=10] node[anchor=south,left]{12} (F2);
				\path[->,>=stealth'] (F2) edge[bend left=10] node[anchor=south,right]{13} (MAIN);
				\path[->,>=stealth'] (MAIN) edge[bend left=10] node[anchor=south,right]{14} (F3);
				\path[->,>=stealth'] (F3) edge[bend left=10] node[anchor=south,right]{15} (MAIN);
			
				\draw[-triangle 45] (MAIN.south west) -- node[anchor=south,left] {1} (VERIF.100);
				\draw[-triangle 45] (VERIF.80) -- node[anchor=south,right] {2} (MAIN.197);
				\draw[-triangle 45] (MAIN.210) |- node[anchor=south,left] {3} (ANALYS.5);
				\draw[-triangle 45] (ANALYS.355) -| node[anchor=south,below] {4} (MAIN.215);			
			
				\draw[-triangle 45] (MAIN.south east) -- node[anchor=south,right] {5} (ADD1.80);
				\draw[-triangle 45] (ADD1.110) -- node[anchor=south,left] {6} (MAIN.343);
				\draw[-triangle 45] (MAIN.335) |- node[anchor=south,right] {7} (ADD2.175);
				\draw[-triangle 45] (ADD2.185) -| node[anchor=south,below] {8} (MAIN.330);				
				\draw[-triangle 45] (MAIN.290) |- node[anchor=north,right] {9} (ADD3.175);
				\draw[-triangle 45] (ADD3.185) -| node[anchor=south,below] {10} (MAIN.270);
			
			\end{scope}
			%Légende
			\begin{scope}
				\node (LEGENDE) at (-7,-5) {\textbf{Légende :}};
				\node (FAMILLE) at (-4.5,-5) [rectangle,draw] {\begin{tabular}{c}Package\end{tabular}};
				\node (MODULE) at (-2,-5) [rectangle,draw,dashed] {\begin{tabular}{c}Module\end{tabular}};
				\path[->,>=stealth'] (0.5,-5.3) edge[bend left=0] node[anchor=south,above]{informations transmises} (3,-5.3);
			\end{scope}
			\end{tikzpicture}
			\caption{Organigramme des différents modules du logiciel}\label{fig:M1}
		\end{figure}
			
		\textbf{Notes :}\\
			(1) Fichier CSV : lancement de la vérification de son format\\
			(2) Code d'erreur : fichier OK ou ERREUR\\
			(3) Fichier CSV : lancement de l'analyse de son contenu\\
			(4) \textbf{structure 1} : contenant les données du fichier, le nombre de lignes et le nombre de colonnes (connus à partir de la taille de la structure)\\
			\hspace*{1.5em} \textbf{structure 2} : contenant 3 informations sur chaque colonne : le type, le rôle et les positions des données erronées ou manquantes\\
			(5) Ensemble de données de type qualitatif\\
			(6) Erreur ou effectifs, effectifs cumulés, fréquences, fréquences cumulées, diagramme en secteur, histogramme\\
			(7) Ensemble de données de type quantitatif discret\\
			(8) Erreur ou indicateurs de tendance central, de dispersion et de forme, les anomalies, la distribution des données, un diagramme à moustaches\\
			(9) Ensemble de données de type quantitatif continu\\
			(10) Même données que (8)\\
			(11) Chemin du fichier CSV importé \\
			(12) Informations du (4) et ensemble de données contenu dans le fichier CSV \\
			(13) Signal de validation du choix de la colonne, et noms des colonnes\\
			(14) Envoi des résultats d'analyses de (6), (8) et (10)\\
			(15) Signal de contrôle : demande d'exportation des résultats de l'ADD, analyse d'une autre colonne, ou importation d'un autre fichier\\
		
		
			%Modifier les fleches aussi ?
			%légende (non modifiée, on la met en accord avec le dernier produit quand on aura fini) ============
			
		\subsection{Format du fichier CSV}
		Le format du fichier a été établi par le client \textbf{\textit{DCbrain}}. Son contenu est assez structuré et est décrit par des colonnes aux types prédéfinies :
		\begin{center}\begin{longtable}{|>{\centering}m{5cm}|>{\centering}m{2cm}|>{\centering}m{2cm}|>{\centering}m{2.5cm}|>{\centering\arraybackslash}m{2cm}|}			
			\hline \multicolumn{1}{|c|}{\textbf{timestamp}} & \multicolumn{1}{c|}{\textbf{parent}} & \multicolumn{1}{ c|}{\textbf{enfant}} & \multicolumn{1}{c|}{\textbf{mesure 1}} & {\textbf{mesure 2}} \\
			\hline 	January 1st 2017, 15:00:00.000 & 102 & 95 & 26644.235 & 176.253\\
			\hline
		\end{longtable}\vspace{-2em}\end{center}
		\textbf{Remarque :} le graphe de flux utilisé par \textbf{\textit{DCbrain}} pour analyser les réseaux de ses clients est orienté, d'où l'utilisation des nœuds \textit{parent - enfant} numérotés. Ceux-ci permettent d'identifier une connexion précise.\\
		Les colonnes \textit{mesure} représentent des données mesurées sur une connexion à un temps donné. Le nombre de ces colonne n'est pas limité et en pratique il ne dépassera pas la dizaine.
				
		
		\subsection{Fonctionnalités des modules}
		%Faites un tri dans vos fonctionnalités : réduire le plus possible, pour que tout rentre en 10 pages. Et mettre que les fonctionnalités du produit final.
		%Cette partie a aussi besoin d'un peu de mise en page, il y a trop de marge à droite a cause des enum + itemize + description)
	
		\begin{description}[style=unboxed,leftmargin=0.2cm]
		
			\item\textbf{1. API 1 - Chargement des données}
			\begin{enumerate}
			
				\item Module Vérification format fichier :\\
				\begin{itemize}
					\item Vérification de la possibilité d'ouverture du fichier.
					\item Vérification du contenu:  le fichier est un CSV contenant du texte brut non formaté. (pas de mise en forme avec des balises ou autres marqueurs)
					\item Vérification de l'accessibilité du fichier en lecture.
				\end{itemize}
			
				\item Module Analyse contenu fichier :\\
				Ce module comprend deux fonctionnalités principales :
				\begin{itemize}
					\item Lecture du contenu du fichier CSV :\\
					On initialise une première structure (\textbf{structure 1}) pour y sauvegarder le contenu du fichier.\\
					On lit ligne par ligne des caractères du fichier. A chaque fois qu'on détecte un caractère de séparation (une virgule, un point-virgule ou une tabulation), on stocke les caractères lus (la donnée) dans la structure.\\
					Cette fonctionnalité fournit la \textbf{structure 1} contenant les données du fichier, le nombre de lignes et le nombre de colonnes (connus à partir de la taille de la structure).\\
					
					\item Descriptions préliminaires des données de chaque colonne du fichier CSV :\\
					On initialise une deuxième structure (\textbf{structure 2}) pour y stocker des informations sur chaque colonne.\\
					On lit une par une les données de chaque colonne. On vérifie le type de la donnée en le comparant au type attendu. Si le type ne correspond pas, on signale dans la structure que la colonne contient une donnée erronée ou manquante (repérée par sa position dans la colonne). A la fin du parcours d'une colonne, on pourra lui attribuer un rôle/nom.\\
					Cette fonctionnalité fournit donc la \textbf{structure 2} contenant 3 informations sur chaque colonne : le type, le rôle et les positions des données erronées ou manquantes.
				\end{itemize}
				Ce module va donc fournir les deux structures \textbf{structure 1} et \textbf{structure 2} décrites ci-dessus.
			\end{enumerate}
			~\\
			\item\textbf{2. API 2 - Analyse descriptive de données}
			\begin{enumerate}
				\item Module ADD qualitatives :
					\begin{itemize}
					\item Calcul des effectifs, effectifs cumulés, fréquences et fréquences cumulées des données
					\item Préparation des représentations graphiques : diagramme en secteur et histogramme
					\end{itemize}
				\item Module ADD quantitatives discrètes :
					\begin{itemize}
					\item Statistiques : moyenne, quantiles, variance, écart-type, coefficients de symétrie et d'aplatissement de Fisher
					\item Anomalies de Tukey : toute donnée hors de l'intervalle [Q1 - 1.5*IQ ; Q3 + 1.5*IQ], où Q1 et Q3 sont les quartiles, et IQ l'écart inter-quartiles.
					\item Préparation des représentations graphiques : fonction de distribution, fonction de distribution cumulative,  boîte à moustaches
					\end{itemize}
				\item Module ADD quantitatives continues :
					\begin{itemize}
					\item Discrétisation de l'étendue (découpage en classe d'intervalles)
					\item Calcul des quantiles pour le cas  de données continues
					\item Préparation des représentations graphiques : fonction de distribution cumulative
					\end{itemize}
			\end{enumerate}
			~\\
			\item\textbf{3. Interface web}
			\begin{enumerate}
				\item Module Gestion des flux :
					\begin{itemize}
					\item Gestion des branchements : exécution normale ou arrêts pour cause d'erreur.
					\item Interface entres les différentes fonctionnalités : communique les données nécessaires entre les modules.
					\end{itemize}
					
				\item Module Fenêtre choix fichier :
					\begin{itemize}
					\item Récupération un fichier CSV en renseignant son chemin en parcourant l'arborescence de fichiers, ou de la manière d'un Drag \& Drop.
					\end{itemize}
					
				\item Module Fenêtre rôle et choix colonne :
					\begin{itemize}
					\item Affichera le nombre de lignes/colonnes contenu dans le CSV.
					\item Affichera le titre du fichier.
					\item Affichera un échantillon du contenu du CSV (environ les 1000 premières lignes) avec un système de scroll.
					\item Affichage des lignes erronés (numéro de la ligne + contenu + type d'erreur).
					\item Mise en place d'un système de navigation sous forme d'onglet (Onglet erreurs, onglet échantillon,...). Cela permettra d'éviter que la fenêtre contienne trop d'informations.
					\item L'Utilisateur devra sélectionner la colonne avec un clic, puis pourra lancer l'analyse sur celle-ci. 
					\end{itemize}
				
				\item Module Fenêtre résultats ADD :
					\begin{itemize}
					\item Affichage des résultats d'analyse descriptive : informations statistiques de l'API 2 et représentations graphiques pour visualiser les données dans leur ensemble.
					\item Une fonctionnalité de retour en arrière permet de sélectionner une nouvelle colonne sans relancer l'applet.
					\end{itemize}
				\end{enumerate}
			
		\end{description}
		
	\section{Contraintes et Langages de programmation}
		\subsection{Contraintes}
		Le développement de l'application s'accompagne de contraintes imposées. Celles-ci établie par DCbrain au commencement du projet. On peut en énumérer trois, qui sont les suivantes:\\
		\begin{center}\begin{longtable}{|>{\centering}m{4cm}|>{\raggedright\arraybackslash}m{11cm}|}			
				\hline \multicolumn{1}{|c}{\textbf{Contrainte}} & \multicolumn{1}{|c|}{\textbf{Fiche}} \\
				\hline 	1. Le produit doit fournir une application web &
						\begin{description}[style=unboxed,leftmargin=0.2cm]
						\item{\textbf{Description :}} notre produit sera une application fonctionnant sur un navigateur web, appelée \textbf{\textit{applet}}.
						\item{\textbf{Justification :}} assure une très grande portabilité et fournit à l'utilisateur une interface interactive.
						\item{\textbf{Critère de satisfaction :}} on peut lancer l'application sur un navigateur web.
						\end{description}\\
				\hline 2. Le produit doit être développé avec un langage de programmation compatible avec l'analyse de données &
						\begin{description}[style=unboxed,leftmargin=0.2cm]
						\item{\textbf{Description :}} le langage de programmation choisi doit inclure des bibliothèques qui permettent d'analyser les données (taches de data mining et de machine Learning).
						\item{\textbf{Justification :}} permettre une analyse de données efficace.
						\item{\textbf{Critère de satisfaction :}} on peut effectuer une analyse descriptive de données.
						\end{description}\\
				\hline 3. Le produit doit fournir une API d'analyse de données en sortie &
						\begin{description}[style=unboxed,leftmargin=0.2cm]
						\item{\textbf{Description :}} l'application doit intégrer des API d'analyse descriptive de données qui pourront être livrées en sortie au client.
						\item{\textbf{Justification :}} permettre une réutilisabilité des fonctionnalités majeures du produit.
						\item{\textbf{Critère de satisfaction :}} on peut exporter une API d'analyse de données en sortie.
						\end{description}\\
				\hline
				\end{longtable}\vspace{1em}\end{center}
			
			
		\subsection{Choix langages et frameworks}
		Les contraintes appliquées au projet ainsi que les exigences définies nous ont permis de fixer notre choix sur un langage de programmation : \textbf{\textit{Python}}.\\
		Ce choix est justifié par le fait que \textbf{\textit{Python}} soit compatible pour l'Analyse Descriptives de Données, puisqu'il est doté de nombreux modules permettant d'effectuer du calcul scientifique sur des données, donc l'analyse de données. De même, il est compatible avec le développement d'applications web, grâce aux nombreux frameworks disponibles.\\
		 Pour ce dernier, notre choix s'est porté sur \textbf{\textit{Flask}}, qui s’efforce d’être le plus simple possible pour une prise en main et le plus petit possible. À l'opposé de Django ou autres, il est donc plus abordable pour une première approche sur le développement d'applications web.\\
		
		L'application devant être une applet intégrée dans une page web pour s'exécuter, le recours à des langages de balisage comme \textbf{\textit{HTML}} pour la présentation et \textbf{\textit{CSS}} pour la mise en forme de ces pages web étaient nécessaires.\\
		
		Ensuite, pour permettre à l'application de faire communiquer les pages web avec un serveur web sans occasionner le rechargement de la page, nous avons eu recours à \textbf{\textit{Ajax}}. Il s'agit d'un concept de programmation web reposant sur plusieurs technologies comme le \textbf{\textit{Javascript}}, qui lui se charge d'établir la connexion entre la page web et serveur. \textbf{\textit{Javascript}} est également employé pour rendre les pages web interactives. \textbf{\textit{Jquery}} est quant à lui utilisé pour faire le lien entre \textbf{\textit{HTML}} et \textbf{\textit{Ajax}}.\\
		
		Enfin, pour la génération de la documentation des fonctionnalités du programme, on utilisera le logiciel \textbf{\textit{Sphinx}} qui se charge de généré la docstring à partir des commentaires présents dans le code python.\\
		
	\section{Partie technique du projet}
	Quoi marche (comment ca marche - à l'aide de la partie avant)
	QUoi marche pas (pb et et resolution si trouvee)
	
	1 page, plus tard
		
	\section{Points délicats}
		
	Justification des points qui divergent par rapport aux specs et organigramme de départ
	
	1 page, plus tard à moins que t'ai déjà des idées 
	- (Peut-être rajouté les points délicats de flask/application web sur l'approche et tout ce qui va avec?
				
	\section{Organisation interne du groupe}
	Ci-dessous se trouve le tableau de la répartition des modules pour chaque membre du groupe établie avant le commencement du développement:
	\begin{center}\begin{longtable}{|>{\centering}m{5cm}|>{\centering}m{2cm}|>{\centering}m{2cm}|>{\centering}m{2.5cm}|>{\centering\arraybackslash}m{1cm}|}			
		\hline \multicolumn{1}{|c|}{\textbf{Module}} & \multicolumn{1}{c|}{\textbf{Malek}} & \multicolumn{1}{ c|}{\textbf{Sonny}} & \multicolumn{1}{c|}{\textbf{Jean-Didier}} & {\textbf{Total}} \\
		\hline 	Gestion des flux & ~ & ~ & x & 1\\
		\hline 	Fenêtre choix fenêtre & ~ & ~ & x & 1\\
		\hline 	Fenêtre rôle et choix colonne & x & ~ & ~ & 1\\
		\hline 	Fenêtre résultats ADD & ~ & x & ~ & 1\\
		\hline  ADD qualitatives & ~ & ~ & x & 1\\
		\hline 	ADD quantitatives discrètes & ~ & x & ~ & 1\\
		\hline 	ADD quantitative continues &  ~ & x & ~ & 1\\
		\hline 	Vérification format fichier & x & ~ & ~ & 1\\
		\hline 	Analyse contenu fichier & x & ~ & ~ & 1\\
		\hline
		\end{longtable}\vspace{1em}\end{center}
		
	1 page
	On expliquera plus tard comment ça s'est déroulé concrètement par rapport à ce qui était prévu.
		
	\section{Coût}
	Tableau indiquant l'estimation et le coût final en nombre de lignes de code, pour chaque module:
	\begin{center}\begin{longtable}{|>{\centering}m{3cm}|>{\centering}m{3cm}|>{\centering\arraybackslash}m{7cm}|>{\centering}m{1.5cm}|}			
		\hline \multicolumn{1}{|c|}{\textbf{Module}} & \multicolumn{1}{c|}{\textbf{Nombre de lignes}} & \multicolumn{1}{c|}{\textbf{Justification}} & \multicolumn{1}{c|}{\textbf{Coût final}}\\
		\hline 	Gestion des flux & 15 & Mise en forme du main et appel de l'application & \\
		\hline 	Fenêtre choix fichier & 10 + 20 & Fonctions pour : Drag\& Drop + Système de fichiers &\\
		\hline 	Fenêtre rôle et choix colonne & 5 + 20 + 10 + 10 + 20 & Communication avec le module application + Affichage de l'interface+ lecteurs des valeurs + 	affichage des valeurs & \\
		\hline 	Fenêtre résultats ADD & 10 + 3* 30 &  Envoie d'informations au module application + Construction des graphe pour l'ADD pour les 3 types d'analyse & \\
		\hline  ADD qualitatives  & 20 + 20*3 & Application des formules pour les calculs de fréquences et d'effectifs + calcul des valeurs pour la construction de 3 graphes & \\
		\hline 	ADD quantitatives discrètes & 60 + 20*2 & Application des formules attaché à l'analyse quantitative discret + calcul des valeurs pour la construction de 2 graphes & \\
		\hline 	ADD quantitatives continu & 20 + 10 + 10 + 5 + 20*2 & Parcours + choix précision classe d'intervalle + écriture + communication avec les modules+ calcul des valeurs pour la construction de 2 graphes & \\
		\hline 	Vérification format fichier & 30 & Ouverture fichier + vérification si ouverture en lecture + présence de texte formaté ou non &\\
		\hline 	Analyse contenu fichier & 20 + 5 + 25 + 10 &  Recopie et vérification + initialisation de la structure+ Parcours du fichiers avec condition + Fonction pour donner nom et type de colonne & \\
		\hline \textbf{Coût Total} & \textbf{565} & \textbf{Estimation totale du coût} & \\
		\hline 	
	\end{longtable}\vspace{1em}\end{center}
		
	\section*{Conclusion}
	1 demi-page, plus tard
		
\end{document}
